\subsection{纤维群胚代数的连续对称：Lie 型选择律与 \texorpdfstring{$\mathrm{U}(1)$}{U(1)} 阻断}\label{subsec:foldbin-groupoid-aut-lie-selection}
在二进制区间折叠 \(\Fold_m^{\mathrm{bin}}\) 的口径下，纤维等价关系自然生成一条有限群胚；其复群胚代数为有限维半单代数，并由退化直方图决定（命题 \ref{prop:fold-bin-groupoid-algebra-wedderburn}）。
本小节抽取其连通连续对称的完全闭式。

\begin{theorem}[连通 \texorpdfstring{$\ast$}{*}-自同构群的 Lie 型选择律]\label{thm:foldbin-groupoid-aut0-pu-product}
\leanverified{paper\_foldbin\_groupoid\_aut0\_pu\_product}
设 \(\mathcal{R}_m^{\mathrm{bin}}\) 为 \(\Fold_m^{\mathrm{bin}}\) 的纤维对群胚，\(\CC[\mathcal{R}_m^{\mathrm{bin}}]\) 为其复群胚代数，并记纤维多重度 \(d_m^{\mathrm{bin}}(w)=|(\Fold_m^{\mathrm{bin}})^{-1}(w)|\)（\(w\in X_m\)）。
则其 \(\ast\)-自同构群的恒等连通分量满足自然同构
\[
\mathrm{Aut}\!\bigl(\CC[\mathcal{R}_m^{\mathrm{bin}}]\bigr)^0
\ \cong\
\prod_{w\in X_m}\mathrm{PU}\!\left(d_m^{\mathrm{bin}}(w)\right),
\]
从而其 Lie 代数满足
\[
\mathfrak{aut}^0
\ \cong\
\bigoplus_{w\in X_m}\mathfrak{su}\!\left(d_m^{\mathrm{bin}}(w)\right).
\]
因此，作为连通连续因子的简单 Lie 型只能来自 \(A_{d-1}\)（即 \(\mathfrak{su}(d)\)），且 \(d\) 必出现在退化谱取值集合 \(\{d_m^{\mathrm{bin}}(w):w\in X_m\}\) 中。
\end{theorem}
\begin{proof}
由命题 \ref{prop:fold-bin-groupoid-algebra-wedderburn}，存在 \(\ast\)-代数同构
\(
\CC[\mathcal{R}_m^{\mathrm{bin}}]\cong \bigoplus_{w\in X_m} M_{d_m^{\mathrm{bin}}(w)}(\CC)
\)。
对每个矩阵块 \(M_d(\CC)\)，其 \(\ast\)-自同构群的连通分量由内自同构给出，等价于 \(\mathrm{PU}(d)\)。
直和代数的 \(\ast\)-自同构还可离散地置换同型矩阵块，但该置换部分属于离散群，故不贡献于恒等连通分量。
因此连通分量为各块内自同构的直积，得到所示同构；Lie 代数为直和 \(\bigoplus_w\mathfrak{su}(d_m^{\mathrm{bin}}(w))\)。证毕。
\end{proof}

\begin{proposition}[基本群分解与挠指数]\label{prop:foldbin-groupoid-aut0-pi1-torsion-exponent}
令
\[
\mathcal G_m:=\mathrm{Aut}\!\bigl(\CC[\mathcal{R}_m^{\mathrm{bin}}]\bigr)^0.
\]
则基本群满足自然同构
\[
\pi_1(\mathcal G_m)\ \cong\ \prod_{w\in X_m}\ZZ/d_m^{\mathrm{bin}}(w)\ZZ,
\]
并且其挠指数（exponent）为
\[
\exp\bigl(\pi_1(\mathcal G_m)\bigr)
=\operatorname{lcm}\{\,d_m^{\mathrm{bin}}(w):\ w\in X_m\,\}.
\]
\end{proposition}
\begin{proof}
由定理 \ref{thm:foldbin-groupoid-aut0-pu-product} 有
\(
\mathcal G_m\cong\prod_{w\in X_m}\mathrm{PU}(d_m^{\mathrm{bin}}(w))
\)。
又对 \(d\ge 1\) 有标准同构 \(\pi_1(\mathrm{PU}(d))\cong \ZZ/d\ZZ\)，并且
\(
\pi_1(\prod_i G_i)\cong\prod_i\pi_1(G_i)
\)，故得第一式。
挠指数为直积中各循环因子的阶的最小公倍数，即得第二式。证毕。
\end{proof}

\begin{corollary}[共同循环识别的充要条件]\label{cor:foldbin-groupoid-aut0-cyclic-injection-criterion}
取整数 \(n\ge 1\)。
则存在群单射 \(\ZZ/n\ZZ\hookrightarrow \pi_1(\mathcal G_m)\) 当且仅当
\[
n\ \mid\ \exp\bigl(\pi_1(\mathcal G_m)\bigr).
\]
\end{corollary}
\begin{proof}
\(\pi_1(\mathcal G_m)\) 为有限阿贝尔群且其指数即其元素阶的最小公倍数。
因此 \(\pi_1(\mathcal G_m)\) 含有阶为 \(n\) 的元素当且仅当 \(n\) 整除其指数；而 \(\ZZ/n\ZZ\) 的群单射等价于存在阶为 \(n\) 的元素。证毕。
\end{proof}

\begin{theorem}[挠素因子谱的无限性排除统一有限维连通李包络]\label{thm:foldbin-groupoid-aut0-no-universal-lie-limit}
\leanverified{paper\_foldbin\_groupoid\_aut0\_no\_universal\_lie\_limit}
令 \(M\subseteq 2\NN\) 为无限集，并令
\[
\mathcal P_M
:=\bigcup_{m\in M}\Bigl\{\,p:\ p\ \text{为素数且}\ p\mid \exp\bigl(\pi_1(\mathcal G_m)\bigr)\Bigr\}.
\]
若 \(\mathcal P_M\) 为无限集合，则不存在连通李群 \(L\) 与一族群单射
\[
\iota_m:\ \pi_1(\mathcal G_m)\hookrightarrow \pi_1(L)
\qquad(m\in M).
\]
\end{theorem}
\begin{proof}
任意连通李群 \(L\) 的 \(\pi_1(L)\) 为有限生成阿贝尔群，从而其挠子子群有限；
特别地，\(\pi_1(L)\) 的挠子部分只包含有限多个素数的挠元。
若存在单射族 \(\iota_m\)，则对每个 \(m\in M\) 与每个 \(p\mid \exp(\pi_1(\mathcal G_m))\)，
\(\pi_1(L)\) 必含有阶为 \(p\) 的元素。
于是 \(\mathcal P_M\) 必为有限集，矛盾。证毕。
\end{proof}

\begin{remark}\label{rem:foldbin-groupoid-aut0-torsion-prime-spectrum}
内部文献 \cite{Internal2026FoldbinGroupoidAutPi1TorsionSpectrum20260302} 给出一条可审计的挠谱生成机制：对任意素数 \(p\)，存在偶数分辨率 \(m\) 使
\[
p\ \mid\ \exp\bigl(\pi_1(\mathcal G_m)\bigr),
\]
等价地，\(\pi_1(\mathcal G_m)\) 含有阶为 \(p\) 的元素。
因此 \(M=2\NN\) 的 \(\mathcal P_M\) 为无限集合；
结合定理 \ref{thm:foldbin-groupoid-aut0-no-universal-lie-limit}，这排除了把 \(\{\mathcal G_m\}_{m\in 2\NN}\) 作为某个固定有限维连通李群之连续极限（在基本群单射意义下）的方案。
\end{remark}

\begin{corollary}[window-\(6\) 的 \texorpdfstring{$\mathrm{U}(1)$}{U(1)} 因子阻断]\label{cor:foldbin6-no-u1-in-aut0-groupoid-algebra}
在 \(m=6\) 的二进制区间折叠口径下，纤维退化直方图为 \(2:8,\,3:4,\,4:9\)（推论 \ref{cor:fold-bin-m6-fiber-hist}）。
因此
\[
\mathrm{Aut}\!\bigl(\CC[\mathcal{R}_6^{\mathrm{bin}}]\bigr)^0
\ \cong\
\mathrm{PU}(2)^8\times \mathrm{PU}(3)^4\times \mathrm{PU}(4)^9,
\]
为半单紧 Lie 群，从而其恒等连通分量不含任何 \(\mathrm{U}(1)\) 直因子。
\end{corollary}
\begin{proof}
由定理 \ref{thm:foldbin-groupoid-aut0-pu-product} 直接代入 \(m=6\) 的退化谱取值与出现次数得到乘积分解。
由于 \(\mathrm{PU}(d)\)（\(d\ge 2\)）为紧致半单群，有限直积仍为半单群，故不含 \(\mathrm{U}(1)\) 直因子。证毕。
\end{proof}

\begin{theorem}[连通自同构群的三阶整上同调秩由纤维数完全决定]\label{thm:foldbin-groupoid-aut0-third-cohomology-rank}
\leanverified{paper\_foldbin\_groupoid\_aut0\_third\_cohomology\_rank}
沿用命题 \ref{prop:foldbin-groupoid-aut0-pi1-torsion-exponent} 的记号，记
\[
\mathcal G_m:=\mathrm{Aut}\!\bigl(\CC[\mathcal{R}_m^{\mathrm{bin}}]\bigr)^0
\cong
\prod_{w\in X_m}\mathrm{PU}\!\left(d_m^{\mathrm{bin}}(w)\right).
\]
则有
\[
H^1(\mathcal G_m;\ZZ)=0,
\qquad
H^2(\mathcal G_m;\ZZ)=0,
\]
并且
\[
\boxed{\;
H^3(\mathcal G_m;\ZZ)
\cong
\bigoplus_{\{w\in X_m:\ d_m^{\mathrm{bin}}(w)\ge 2\}}\ZZ.\;}
\]
特别地，若所有纤维大小都至少为 \(2\)，则
\[
\operatorname{rank}H^3(\mathcal G_m;\ZZ)=|X_m|.
\]
在 \(m=6\) 的二进制区间折叠口径下，由推论 \ref{cor:fold-bin-m6-fiber-hist} 的直方图
\[
2:8,\qquad 3:4,\qquad 4:9
\]
可知全部纤维都满足 \(d_6^{\mathrm{bin}}(w)\ge 2\)，且 \(8+4+9=21\)。
因此
\[
H^3\!\Bigl(\mathrm{Aut}\!\bigl(\CC[\mathcal{R}_6^{\mathrm{bin}}]\bigr)^0;\ZZ\Bigr)
\cong
\ZZ^{21}.
\]
\end{theorem}
\begin{proof}
由定理 \ref{thm:foldbin-groupoid-aut0-pu-product}，
\[
\mathcal G_m\cong \prod_{w\in X_m}\mathrm{PU}\!\left(d_m^{\mathrm{bin}}(w)\right).
\]
对 \(d=1\) 有 \(\mathrm{PU}(1)=\{1\}\)。
对每个 \(d\ge 2\)，标准拓扑事实给出
\[
H^1(\mathrm{PU}(d);\ZZ)=0,
\qquad
H^2(\mathrm{PU}(d);\ZZ)=0,
\qquad
H^3(\mathrm{PU}(d);\ZZ)\cong \ZZ.
\]
于是由 K\"unneth 定理，有限直积的 \(H^1\) 与 \(H^2\) 仍为零，而 \(H^3\) 为各因子三阶上同调的直和：
\[
H^3\!\Bigl(\prod_{w\in X_m}\mathrm{PU}\!\left(d_m^{\mathrm{bin}}(w)\right);\ZZ\Bigr)
\cong
\bigoplus_{\{w:\ d_m^{\mathrm{bin}}(w)\ge 2\}}
H^3\!\left(\mathrm{PU}\!\left(d_m^{\mathrm{bin}}(w)\right);\ZZ\right).
\]
代入 \(H^3(\mathrm{PU}(d);\ZZ)\cong\ZZ\) 即得所示同构。
若所有纤维大小都至少为 \(2\)，则每个 \(w\in X_m\) 都贡献一个 \(\ZZ\)-因子，故秩等于 \(|X_m|\)。
在 \(m=6\) 时，由直方图 \(2:8,3:4,4:9\) 的总数为 \(21\)，得到
\[
H^3\!\Bigl(\mathrm{Aut}\!\bigl(\CC[\mathcal{R}_6^{\mathrm{bin}}]\bigr)^0;\ZZ\Bigr)\cong \ZZ^{21}.
\]
证毕。
\end{proof}

\begin{proposition}[window-\(6\) 边界二层覆盖诱导的复 Clifford 子代数]\label{prop:foldbin6-boundary-clifford-subalgebra}
取 \(m=6\) 并令
\(
X^{\mathrm{bdry}}_6:=\{w\in X_6:\ w_1=w_6=1\}
\)。
则每个 \(w\in X^{\mathrm{bdry}}_6\) 在二进制区间折叠 \(\Fold^{\mathrm{bin}}_6\) 下恰有两条预像，且差值严格为 \(34=F_9\)：
\[
(\Fold^{\mathrm{bin}}_6)^{-1}(w)=\{V_6(w),\,V_6(w)+34\}
\qquad(\text{见表 }\ref{tab:fold6_boundary_sheet_pairs}).
\]
在 Wedderburn 分解
\(
\CC[\mathcal{R}_6^{\mathrm{bin}}]\cong \bigoplus_{u\in X_6} M_{d^{\mathrm{bin}}_6(u)}(\CC)
\)
中，边界词 \(w\in X^{\mathrm{bdry}}_6\) 对应的矩阵块为 \(M_2(\CC)\)。
以 sheet parity 诱导的有序基 \((V_6(w),\,V_6(w)+34)\) 识别该块为 \(M_2(\CC)\) 后，定义
\[
P_w:=\begin{pmatrix}1&0\\0&-1\end{pmatrix},
\qquad
X_w:=\begin{pmatrix}0&1\\1&0\end{pmatrix}.
\]
则 \(P_w^2=X_w^2=I\) 且 \(P_wX_w=-X_wP_w\)，从而由 \((P_w,X_w)\) 生成的 \(\ast\)-子代数同构于复 Clifford 代数 \(\mathrm{Cl}_2(\CC)\cong M_2(\CC)\)，并且其 \(\ZZ_2\)-分级与 sheet parity 一致。
对全部边界词取直和，得到一个由边界二层覆盖数据规范确定的 \(\ZZ_2\)-分级子代数
\[
\bigoplus_{w\in X^{\mathrm{bdry}}_6}\mathrm{Cl}_2(\CC)
\ \subset\
\CC[\mathcal{R}_6^{\mathrm{bin}}].
\]
\end{proposition}
\begin{proof}
由推论 \ref{cor:terminal-delta-parity-rule} 与表 \ref{tab:fold6_boundary_sheet_pairs}，边界预像对严格为两点且具自然 sheet parity 标号，故相应矩阵块为 \(M_2(\CC)\) 且可取所述有序基。
在该基下，\(P_w\) 为 parity 算子、\(X_w\) 为交换两 sheet 的矩阵表示，直接计算得 \(P_w^2=X_w^2=I\) 与 \(P_wX_w=-X_wP_w\)。
该关系给出 \(\mathrm{Cl}_2(\CC)\) 的标准生成刻画，而 \(\mathrm{Cl}_2(\CC)\) 与 \(M_2(\CC)\) 的同构为复 Clifford 代数的基本事实。
逐个 \(w\) 的构造彼此独立，故直和嵌入到 Wedderburn 直和中，从而得到所述子代数。证毕。
\end{proof}

\begin{theorem}[二层覆盖的 Clifford 边界代数：由纤维群胚块内生生成的极大 \texorpdfstring{$\ZZ_2$}{Z2}-分级]\label{thm:foldbin6-boundary-clifford-maximal-graded-subalgebra}
沿用命题 \ref{prop:foldbin6-boundary-clifford-subalgebra} 的记号，并记
\[
\mathcal A_{\partial}
:=
C^\ast\bigl(\{P_w,X_w:\ w\in X_6^{\mathrm{bdry}}\}\bigr)
\subset
\CC[\mathcal R_6^{\mathrm{bin}}].
\]
则有规范同构
\[
\mathcal A_{\partial}
=
\bigoplus_{w\in X_6^{\mathrm{bdry}}} M_2(\CC)
\cong
\bigoplus_{w\in X_6^{\mathrm{bdry}}}\mathrm{Cl}_2(\CC),
\]
并且其 \(\ZZ_2\)-分级正由 sheet parity 给出：偶部由各块中的对角子代数组成，奇部由交换两张 sheet 的反对角元组成。

进一步，在“支撑仅落在这些 boundary 二层块上，且尊重上述 sheet-parity 分级”的范畴内，\(\mathcal A_{\partial}\) 是极大分级 \(\ast\)-子代数：任一包含全部 \(P_w\) 与 \(X_w\) 的分级 \(\ast\)-子代数都必等于 \(\mathcal A_{\partial}\)。
\leanverified{paper\_foldbin6\_boundary\_clifford\_maximal\_graded\_subalgebra}
\end{theorem}
\begin{proof}
命题 \ref{prop:foldbin6-boundary-clifford-subalgebra} 已表明：对每个 \(w\in X_6^{\mathrm{bdry}}\)，由 \(P_w\) 与 \(X_w\) 生成的 \(\ast\)-子代数就是对应的 \(M_2(\CC)\) 块，并携带与 sheet parity 一致的 \(\ZZ_2\)-分级。由于不同 \(w\) 的边界二层块在 Wedderburn 分解中彼此独立，直接求直和即得
\[
\mathcal A_{\partial}
=
\bigoplus_{w\in X_6^{\mathrm{bdry}}}M_2(\CC).
\]

为证极大性，只需逐块检验。固定 \(w\)，在对应的 \(M_2(\CC)\) 中令
\[
e_{11}^{(w)}:=\frac{I+P_w}{2},
\qquad
e_{22}^{(w)}:=\frac{I-P_w}{2}.
\]
则 \(e_{11}^{(w)},e_{22}^{(w)}\) 是两条 sheet 的最小投影。再由 \(X_w\) 得到
\[
e_{12}^{(w)}:=e_{11}^{(w)}X_w e_{22}^{(w)},
\qquad
e_{21}^{(w)}:=e_{22}^{(w)}X_w e_{11}^{(w)},
\]
从而四个矩阵单位 \(e_{ij}^{(w)}\) 全部都落在任何同时包含 \(P_w,X_w\) 的 \(\ast\)-子代数中。因此该子代数不可能是真子代数，只能等于整块 \(M_2(\CC)\)。

对全部 \(w\in X_6^{\mathrm{bdry}}\) 逐块相加，便知任何支撑在 boundary 二层块上、尊重 parity 分级且包含全部 \(P_w,X_w\) 的分级 \(\ast\)-子代数，必等于 \(\mathcal A_{\partial}\)。证毕。
\end{proof}

\endinput
